\documentclass[aspectratio=169,11pt]{beamer}
\input{../common.tex}

\title{Logistic Regression}
\subtitle{TTK4260: Multivariate Data Analysis \& Machine Learning}
\author{Adil Rasheed}
\institute{Department of Engineering Cybernetics, NTNU}
\date{Spring 2026}

\begin{document}

% ============================================================
% TITLE SLIDE
% ============================================================
\NTNUCover

% ============================================================
% OUTLINE
% ============================================================
\begin{frame}{Lecture Outline}
\begin{columns}[T]
\column{0.55\textwidth}
\tableofcontents
\column{0.4\textwidth}
\begin{tikzpicture}[scale=0.7]
  % Sigmoid teaser
  \draw[->,thick,BrandBlue!60] (-3,0) -- (3.5,0) node[right,font=\scriptsize]{$z$};
  \draw[->,thick,BrandBlue!60] (0,-0.3) -- (0,2.8) node[above,font=\scriptsize]{$\sigma(z)$};
  \draw[dashed,BrandBlue!30] (-3,2.2) -- (3.3,2.2) node[right,font=\tiny]{$1$};
  \draw[dashed,BrandBlue!30] (-3,1.1) -- (3.3,1.1) node[right,font=\tiny]{$0.5$};
  \draw[very thick,BrandBlue] plot[domain=-3:3,samples=80] (\x, {2.2/(1+exp(-1.8*\x))});
  \node[font=\footnotesize\itshape,BrandBlue] at (0,-0.9) {Preview: the sigmoid function};
\end{tikzpicture}
\end{columns}
\end{frame}

% ============================================================
\section{Motivation \& Problem Setup}
% ============================================================
\begin{frame}{Motivation: Is This Sensor Reading a Fault?}
\begin{columns}[T]
\column{0.50\textwidth}
\begin{center}
\begin{tikzpicture}[scale=0.85]
  % Axes
  \draw[->,thick] (-0.3,0) -- (6.5,0) node[right,font=\small]{pressure (bar)};
  \draw[->,thick] (0,-0.3) -- (0,5.5) node[above,font=\small]{temperature (\textdegree C)};

  % Axis ticks
  \foreach \x/\l in {1/10, 2/20, 3/30, 4/40, 5/50, 6/60}{
    \draw (\x,0.07) -- (\x,-0.07) node[below,font=\tiny]{\l};
  }
  \foreach \y/\l in {1/50, 2/100, 3/150, 4/200, 5/250}{
    \draw (0.07,\y) -- (-0.07,\y) node[left,font=\tiny]{\l};
  }

  % Normal cluster (blue, lower-left)
  \foreach \x/\y in {1.0/1.0, 1.3/1.4, 1.1/0.8, 1.5/1.2, 0.8/1.1,
                     1.4/0.9, 1.7/1.5, 1.2/1.6, 0.9/0.7, 1.6/1.0,
                     1.8/1.3, 1.0/1.3, 1.5/0.7, 1.3/1.7, 0.7/1.0,
                     2.0/1.4, 1.1/1.5, 1.6/1.8, 2.1/1.1, 1.8/0.8}{
    \fill[BrandBlue!80] (\x,\y) circle (2.5pt);
  }

  % Fault cluster (red, upper-right)
  \foreach \x/\y in {4.2/3.8, 4.5/4.2, 4.0/3.5, 4.8/4.0, 4.3/4.5,
                     4.7/3.7, 5.0/4.3, 4.1/4.1, 4.6/3.6, 5.2/4.5,
                     4.4/3.9, 4.9/4.6, 5.1/3.8, 4.3/4.3, 5.3/4.1,
                     4.8/4.4, 4.0/4.0, 5.0/3.5, 4.6/4.7, 4.2/3.6}{
    \fill[red!70] (\x,\y) circle (2.5pt);
  }

  % Decision boundary
  % \draw[very thick,accentorange,dashed] (0.5,5.2) -- (3.8,2.5);

  % Boundary label — in the empty upper-left side of the line
  % \node[font=\scriptsize,accentorange,rotate=-45,
  %       fill=white,inner sep=1pt] at (2.9,4.5) 
  %   {$\boldsymbol{\theta}^\top\boldsymbol{x}+\theta_0=0$};

  % Query point
  \fill[accentorange] (3.0,3.0) circle (4pt);
  \node[font=\small\bfseries,accentorange,left] at (2.85,3.6) {$p=\,$?};

  % Cluster labels
  \node[font=\small\bfseries,BrandBlue] at (1.4,2.3) {normal};
  \node[font=\small\bfseries,red!70] at (4.7,5.2) {fault};
\end{tikzpicture}
\end{center}

\column{0.46\textwidth}
\textbf{The scenario:}\\
A process plant logs pressure and temperature from a reactor vessel. Historical data is labelled \textit{normal} or \textit{fault}.

\vspace{0.5em}
When we plot these two features against each other, \textbf{two clusters} emerge.

\vspace{0.5em}
A new reading arrives at the orange point. We need to answer:

\vspace{0.3em}
\begin{center}
\begin{tikzpicture}
  \node[draw=BrandBlue, fill=SoftBlue, rounded corners=4pt,
        inner sep=6pt, font=\normalsize\bfseries, align=center] {%
    Normal or fault?\\[0.2em]
    \normalfont\normalsize And \textbf{how confident} are we?%
  };
\end{tikzpicture}
\end{center}

\end{columns}
\end{frame}

\begin{frame}{From Regression to Classification}
\begin{columns}[T]
\column{0.55\textwidth}
\textbf{What you already know:}
\begin{itemize}
  \item Linear regression predicts \textbf{continuous} values $\hat{y}\in\mathbb{R}$
  \item MLE with Gaussian noise $\Rightarrow$ MSE loss
  \item Closed-form solution: $\hat{\boldsymbol{\theta}}=(\mathbf{X}^\top\mathbf{X})^{-1}\mathbf{X}^\top\boldsymbol{y}$
\end{itemize}
\vspace{0.6em}
\textbf{But many real problems are categorical:}
\begin{itemize}
  \item Spam vs.\ not spam
  \item Healthy vs.\ diseased tissue
  \item Ship collision risk: yes or no?
  \item Sensor fault detection
\end{itemize}

\column{0.4\textwidth}
\begin{block}{Key Question}
\centering
Can we adapt the linear model\\[0.3em]
$z = \boldsymbol{\theta}^\top\mathbf{X} + \theta_0$\\[0.3em]
so that its output represents\\
a \textbf{probability} $\in[0,1]$\,?
\end{block}
\vspace{0.5em}
\begin{tikzpicture}[scale=0.65]
  \fill[BrandBlue!8] (-2.2,-1.5) rectangle (0,1.5);
  \fill[accentgreen!8] (0,-1.5) rectangle (2.2,1.5);
  \draw[dashed,BrandBlue,thick] (0,-1.5) -- (0,1.5);
  % class 0 dots
  \foreach \x/\y in {-1.6/0.8, -0.9/0.3, -1.3/-0.5, -0.7/1.1, -1.8/-0.9, -0.5/-0.3, -1.1/0.0}{
    \fill[red!70] (\x,\y) circle (2.5pt);
  }
  % class 1 dots
  \foreach \x/\y in {1.5/0.7, 0.8/0.2, 1.2/-0.6, 0.6/1.0, 1.7/-0.8, 0.4/-0.2, 1.0/0.5}{
    \fill[BrandBlue!80] (\x,\y) circle (2.5pt);
  }
  \node[font=\tiny,red!70] at (-1.3,-1.8) {Class 0};
  \node[font=\tiny,BrandBlue] at (1.3,-1.8) {Class 1};
\end{tikzpicture}
\end{columns}
\end{frame}

% ------------------------------------------------------------
\begin{frame}{Why Not Linear Regression for Classification?}
\begin{columns}[T]
\column{0.32\textwidth}
\begin{alertblock}{Unbounded Output}
Linear model gives\\
$\hat{y}\in(-\infty,+\infty)$\\[0.3em]
We need $\hat{y}\in[0,1]$\\
for valid probabilities.
\end{alertblock}

\column{0.32\textwidth}
\begin{alertblock}{Outlier Sensitivity}
A single far-away point\\
can shift the decision\\
boundary dramatically.
\end{alertblock}

\column{0.32\textwidth}
\begin{alertblock}{Wrong Loss Surface}
MSE with a sigmoid\\
is \textbf{non-convex}---\\
bad local minima with\\
gradient descent.
\end{alertblock}
\end{columns}

\vspace{1em}
\begin{center}
\begin{tikzpicture}
  \node[draw=BrandBlue, fill=BrandBlue, text=white, rounded corners=4pt, 
        inner sep=8pt, font=\large\bfseries] {%
    Solution: apply a nonlinear transformation $\sigma(z)$ that squashes $\mathbb{R}\to(0,1)$%
  };
\end{tikzpicture}
\end{center}
\end{frame}

% ============================================================
\section{The Sigmoid Function}
% ============================================================

\begin{frame}{The Logistic (Sigmoid) Function}
\begin{columns}[T]
\column{0.45\textwidth}
\begin{block}{Definition}
\vspace{-0.5em}
\[
  \sigma(z) \;=\; \frac{1}{1+e^{-z}}\,,\qquad z = \boldsymbol{\theta}^\top\boldsymbol{x}+\theta_0
\]
\end{block}
\vspace{0.3em}
\textbf{Key properties:}
\begin{itemize}
  \item $\sigma(z)\to 0$ as $z\to -\infty$
  \item $\sigma(z)\to 1$ as $z\to +\infty$
  \item $\sigma(0)=0.5$ \;\textit{(decision threshold)}
  \item Monotonically increasing, differentiable
  \item \textbf{Elegant derivative:}
  \[
    \sigma'(z) = \sigma(z)\bigl(1-\sigma(z)\bigr)
  \]
\end{itemize}

\column{0.52\textwidth}
\begin{tikzpicture}[scale=0.85]
  % Axes
  \draw[->,thick] (-4,0) -- (4.5,0) node[right,font=\small]{$z$};
  \draw[->,thick] (0,-0.4) -- (0,4.2) node[above,font=\small]{$\sigma(z)$};
  % Grid
  \draw[dashed,gray!40] (-4,3.5) -- (4.3,3.5);
  \draw[dashed,gray!40] (-4,1.75) -- (4.3,1.75);
  % Labels
  % Shaded regions
  \fill[red!6] (-4,0) rectangle (0,3.5);
  \fill[BrandBlue!6] (0,0) rectangle (4.3,3.5);
  \node[font=\footnotesize,red!70] at (-2.5,3.9) {\textbf{Class 0}};
  \node[font=\footnotesize,BrandBlue] at (2.5,3.9) {\textbf{Class 1}};
  % Sigmoid curve
  \draw[very thick,BrandBlue] plot[domain=-4:4,samples=100] (\x, {3.5/(1+exp(-1.5*\x))});
  % Point at z=0
  \fill[BrandBlue] (0,1.75) circle (3pt);
  \draw[->,BrandBlue!60,thick] (1.2,1.0) -- (0.15,1.65);
  \node[font=\scriptsize,BrandBlue] at (2.2,0.8) {threshold};
  \node[left,font=\footnotesize] at (-0.1,3.75) {$1$};
  \node[left,font=\footnotesize] at (-0.1,1.75) {$0.5$};
  \node[left,font=\footnotesize] at (-0.1,-0.50) {$0$};
\end{tikzpicture}
\end{columns}
\end{frame}

% ------------------------------------------------------------
\begin{frame}{From Odds to Logits}
\begin{center}
\begin{tikzpicture}[
    box/.style={draw=BrandBlue, fill=BrandBlue!5, rounded corners=3pt, 
                minimum height=2cm, minimum width=3.5cm, align=center, font=\small},
    arr/.style={->,>=stealth,thick,BrandBlue}
  ]
  \node[box] (prob) at (0,0) {\textbf{Probability}\\[0.3em] $p = P(y{=}1\mid\boldsymbol{x})$};
  \node[box] (odds) at (5,0) {\textbf{Odds}\\[0.3em] $\displaystyle\frac{p}{1-p}$};
  \node[box] (logit) at (10,0) {\textbf{Log-Odds (Logit)}\\[0.3em] $\displaystyle\log\!\frac{p}{1-p}=\boldsymbol{\theta}^\top\!\boldsymbol{x}+\theta_0$};
  \draw[arr] (prob) -- (odds);
  \draw[arr] (odds) -- (logit);
\end{tikzpicture}
\end{center}

\vspace{0.5em}
\begin{block}{The Logistic Regression Assumption}
We assume the \textbf{log-odds} of the positive class is a \textbf{linear function} of the features:
\[
  \log\frac{p}{1-p} \;=\; \theta_1 x_1 + \theta_2 x_2 + \cdots + \theta_p x_p + \theta_0
\]
Inverting the logit gives us the sigmoid: \quad $p = \sigma(\boldsymbol{\theta}^\top\boldsymbol{x}+\theta_0)$
\end{block}

\smallskip
\footnotesize\textit{Analogy: linear regression assumes $y = \boldsymbol{\theta}^\top\boldsymbol{x}+\theta_0+\varepsilon$; logistic regression assumes a linear relationship with the \textbf{log-odds}, not $y$ directly. This is a Generalized Linear Model (GLM).}
\end{frame}

% ------------------------------------------------------------
\begin{frame}{Notation Convention}
\begin{center}
\begin{tikzpicture}
  % Absorbed-bias notation
  \node[draw=BrandBlue, fill=SoftBlue, rounded corners=5pt, inner sep=10pt, font=\normalsize, align=left] (absorb) {%
    \textbf{Absorbed-bias notation:} \; define $\tilde{\boldsymbol{x}} = [1,\, x_1,\, x_2,\,\ldots,\, x_p]^\top$ and
    $\tilde{\boldsymbol{\theta}} = [\theta_0,\, \theta_1,\,\ldots,\,\theta_p]^\top$ \\[0.6em]
    Then: \; $z = \boldsymbol{\theta}^\top\boldsymbol{x} + \theta_0 \;=\; \tilde{\boldsymbol{\theta}}^{\,\top}\tilde{\boldsymbol{x}}$
    \hspace{2em}\textit{(a single dot product)}
  };
\end{tikzpicture}
\end{center}

\vspace{0.5em}
\begin{columns}[T]
\column{0.48\textwidth}
\begin{block}{Symbol Conventions}
\begin{tabular}{lll}
  \textbf{Entity} & \textbf{Style} & \textbf{Example} \\[0.3em]
  Scalar & italic, lowercase & $z,\; y,\; \theta_j$ \\
  Vector & bold, italic, lowercase & $\boldsymbol{x},\; \boldsymbol{\theta},\; \boldsymbol{y}$ \\
  Matrix & bold, upright, uppercase & $\mathbf{X},\; \mathbf{H}$
\end{tabular}
\end{block}

\column{0.48\textwidth}
\begin{block}{Key Symbols}
\begin{tabular}{ll}
  $\boldsymbol{\theta}$ & weight vector (learnable) \\
  $\theta_0$            & bias / intercept (learnable) \\
  $\boldsymbol{x}$      & feature vector (input) \\
  $\mathbf{X}$          & design matrix ($n \times p$) \\
  $\boldsymbol{y}$      & target vector \\
  $\sigma(\cdot)$       & sigmoid function
\end{tabular}
\end{block}
\end{columns}

\smallskip
\centering\footnotesize\textit{Throughout these slides we write $\boldsymbol{\theta}$, $\theta_0$ for all learnable parameters; the absorbed form $\tilde{\boldsymbol{\theta}}$ is used where convenient.}
\end{frame}

% ------------------------------------------------------------
\begin{frame}{The Logistic Regression Model}
\begin{center}
\begin{tikzpicture}
  % Main formula box
  \node[draw=BrandBlue, fill=SoftBlue, rounded corners=5pt, inner sep=12pt, font=\Large] (main) {%
    $P(y{=}1 \mid \boldsymbol{x};\,\boldsymbol{\theta},\theta_0) \;=\; \sigma(\boldsymbol{\theta}^\top\boldsymbol{x}+\theta_0) \;=\; \dfrac{1}{1+e^{-(\boldsymbol{\theta}^\top\boldsymbol{x}+\theta_0)}}$%
  };
  
  \node[below=0.5cm of main, draw=BrandBlue!60, fill=white, rounded corners=3pt, inner sep=8pt, font=\normalsize] (comp) {%
    $P(y{=}0 \mid \boldsymbol{x};\,\boldsymbol{\theta},\theta_0) \;=\; 1 - \sigma(\boldsymbol{\theta}^\top\boldsymbol{x}+\theta_0)$%
  };
\end{tikzpicture}
\end{center}

\vspace{0.5em}
\begin{block}{Compact Bernoulli Form}
\[
  P(y\mid\boldsymbol{x};\,\boldsymbol{\theta},\theta_0) \;=\; \sigma(z)^{\,y}\;\bigl(1-\sigma(z)\bigr)^{1-y} \qquad\text{where }z = \boldsymbol{\theta}^\top\boldsymbol{x}+\theta_0
\]
\end{block}
\smallskip
\begin{itemize}
  \item When $y=1$: $P = \sigma(z)$ \quad (first factor survives)
  \item When $y=0$: $P = 1-\sigma(z)$ \quad (second factor survives)
\end{itemize}
\smallskip
\centering\footnotesize\textit{This compact form is essential for the likelihood derivation.}
\end{frame}

% ============================================================
\section{Maximum Likelihood Estimate and The Cost Function}
% ============================================================

\begin{frame}{Maximum Likelihood Estimation}
\textit{You already know MLE---now we apply it to classification.}

\vspace{0.5em}
\textbf{Step 1: Likelihood} (assuming i.i.d.\ observations)
\[
  L(\boldsymbol{\theta},\theta_0) = \prod_{i=1}^{n} P(y_i\mid\boldsymbol{x}_i;\,\boldsymbol{\theta},\theta_0)
               = \prod_{i=1}^{n} \sigma(z_i)^{y_i}\bigl(1-\sigma(z_i)\bigr)^{1-y_i}
\]

\textbf{Step 2: Log-likelihood}
\[
  \ell(\boldsymbol{\theta},\theta_0) = \sum_{i=1}^{n}\Bigl[y_i\log\sigma(z_i) + (1-y_i)\log\bigl(1-\sigma(z_i)\bigr)\Bigr]
\]

\textbf{Step 3: Negative log-likelihood $\to$ our loss function (minimize!)}

\[
  \boxed{J(\boldsymbol{\theta},\theta_0) = -\frac{1}{n}\sum_{i=1}^{n}\Bigl[y_i\log\sigma(z_i) + (1-y_i)\log\bigl(1-\sigma(z_i)\bigr)\Bigr]}
\]

\smallskip
\footnotesize\textit{Compare: MLE with Gaussian noise $\Rightarrow$ MSE (linear regression). MLE with Bernoulli $\Rightarrow$ cross-entropy (logistic regression). Same principle, different distributions!}
\end{frame}

% ------------------------------------------------------------
\begin{frame}{Binary Cross-Entropy Loss: Intuition}
\vspace{-1.5em}
\begin{columns}[T]
\column{0.48\textwidth}
\begin{center}
\textbf{When $y=1$: \; cost $= -\log(p)$}
\begin{tikzpicture}[scale=0.75]
  \draw[->,thick] (0,0) -- (4.5,0) node[right,font=\small]{$p$};
  \draw[->,thick] (0,0) -- (0,4.5) node[above,font=\small]{cost};
  \draw[very thick,BrandBlue] plot[domain=0.05:4,samples=80] ({(\x)}, {-ln(\x/4)*1.3});
  \node[font=\scriptsize,red,right] at (0.3,3.8) {High cost when};
  \node[font=\scriptsize,red,right] at (0.3,3.3) {$p\to 0$ (wrong!)};
  \node[font=\scriptsize,accentgreen,right] at (3.0,1.0) {Low cost};
  \node[font=\scriptsize,accentgreen,right] at (3.0,0.5) {when $p\to 1$};
  \node[left,font=\tiny] at (0,0) {$0$};
  \node[below,font=\tiny] at (4,0) {$1$};
\end{tikzpicture}
\end{center}

\column{0.48\textwidth}
\begin{center}
\textbf{When $y=0$: \; cost $= -\log(1-p)$}
\begin{tikzpicture}[scale=0.75]
  \draw[->,thick] (0,0) -- (4.5,0) node[right,font=\small]{$p$};
  \draw[->,thick] (0,0) -- (0,4.5) node[above,font=\small]{cost};
  \draw[very thick,red!70] plot[domain=0.05:4,samples=80] ({(\x)}, {-ln(1-\x/4.2)*1.3});
  \node[font=\scriptsize,red!70,left] at (4.0,3.8) {High cost when};
  \node[font=\scriptsize,red!70,left] at (4.0,3.3) {$p\to 1$ (wrong!)};
  \node[font=\scriptsize,accentgreen] at (1.0,1.0) {Low cost};
  \node[font=\scriptsize,accentgreen] at (1.0,0.5) {when $p\to 0$};
  \node[left,font=\tiny] at (0,0) {$0$};
  \node[below,font=\tiny] at (4,0) {$1$};
\end{tikzpicture}
\end{center}
\end{columns}

\begin{exampleblock}{Key Property}
\begin{itemize}
  \item The loss penalizes \textbf{confident wrong predictions} severely (cost $\to\infty$)
  \item The cross-entropy loss is \textbf{convex} in $(\boldsymbol{\theta},\theta_0)$ --- guaranteed global optimum!
\end{itemize}
\end{exampleblock}
\end{frame}

% ============================================================
\section{Gradient Descent and Optimization}
% ============================================================

\begin{frame}{Gradient of the Loss Function}
Using the chain rule and $\sigma'(z)=\sigma(z)(1-\sigma(z))$:

\vspace{0.5em}
\textbf{Per-parameter gradient:}
\[
  \frac{\partial J}{\partial \theta_j} = \frac{1}{n}\sum_{i=1}^{m}\bigl(\sigma(z_i)-y_i\bigr)\,x_{ij}
  \qquad\qquad
  \frac{\partial J}{\partial \theta_0} = \frac{1}{n}\sum_{i=1}^{m}\bigl(\sigma(z_i)-y_i\bigr)
\]

\textbf{Vector form:}
\begin{block}{}
\vspace{-0.8em}
\[
  \boxed{\nabla_{\boldsymbol{\theta}} J = \frac{1}{m}\,\mathbf{X}^\top\!\bigl(\sigma(\mathbf{X}\boldsymbol{\theta}+\theta_0\,\boldsymbol{1})-\boldsymbol{y}\bigr)}
\]
\vspace{-0.8em}
\end{block}

\vspace{0.5em}
\begin{exampleblock}{Remarkable Similarity to Linear Regression}
\begin{tabular}{ll}
  Linear regression: & $\nabla_{\boldsymbol{\theta}} J = \frac{1}{m}\,\mathbf{X}^\top(\mathbf{X}\boldsymbol{\theta}-\boldsymbol{y})$ \\[0.5em]
  Logistic regression: & $\nabla_{\boldsymbol{\theta}} J = \frac{1}{m}\,\mathbf{X}^\top\bigl(\sigma(\mathbf{X}\boldsymbol{\theta}+\theta_0\,\boldsymbol{1})-\boldsymbol{y}\bigr)$
\end{tabular}
\smallskip\\
\textit{Same form---the only difference is $\sigma(\cdot)$ applied to the linear prediction!}
\end{exampleblock}
\end{frame}

% ------------------------------------------------------------
\begin{frame}{Gradient Derivation (Step by Step)}
\vspace{0.3em}
\textbf{Start from the loss for a single sample} ($z_i = \boldsymbol{\theta}^\top\boldsymbol{x}_i + \theta_0$):
\[
  J_i = -\bigl[y_i\log\sigma(z_i) + (1-y_i)\log(1-\sigma(z_i))\bigr]
\]

\textbf{Apply the chain rule:} \; $\dfrac{\partial J_i}{\partial \theta_j} = \dfrac{\partial J_i}{\partial \sigma}\cdot\dfrac{\partial \sigma}{\partial z_i}\cdot\dfrac{\partial z_i}{\partial \theta_j}$

\vspace{0.3em}
\begin{columns}[T]
\column{0.30\textwidth}
\begin{block}{Factor 1}
\vspace{-0.5em}
\[
  \frac{\partial J_i}{\partial \sigma} = -\frac{y_i}{\sigma}+\frac{1-y_i}{1-\sigma}
\]
\end{block}

\column{0.30\textwidth}
\begin{block}{Factor 2}
\vspace{-0.5em}
\[
  \frac{\partial \sigma}{\partial z_i} = \sigma(1-\sigma)
\]
\end{block}

\column{0.30\textwidth}
\begin{block}{Factor 3}
\vspace{-0.5em}
\[
  \frac{\partial z_i}{\partial \theta_j} = x_{ij}
\]
\end{block}
\end{columns}

\vspace{0.5em}
\textbf{Multiply and simplify} (the $\sigma(1{-}\sigma)$ cancels elegantly):
\[
  \frac{\partial J_i}{\partial \theta_j}
  = \bigl(\sigma(z_i)-y_i\bigr)\,x_{ij}
  \qquad\Longrightarrow\qquad
  \nabla_{\boldsymbol{\theta}}J = \frac{1}{n}\sum_{i=1}^{n}(\sigma(z_i)-y_i)\,\boldsymbol{x}_i
  = \frac{1}{m}\,\mathbf{X}^\top(\hat{\boldsymbol{p}}-\boldsymbol{y})
\]
\end{frame}

% ------------------------------------------------------------
\begin{frame}{Gradient Descent Algorithm}
\begin{columns}[T]
\column{0.52\textwidth}
\begin{block}{Algorithm}
\begin{enumerate}
  \item Initialize $\boldsymbol{\theta}=\boldsymbol{0}$, $\theta_0=0$
  \item \textbf{Repeat} until convergence:
  \begin{enumerate}
    \item Compute predictions: $\hat{\boldsymbol{p}} = \sigma(\mathbf{X}\boldsymbol{\theta}+\theta_0\,\boldsymbol{1})$
    \item Compute gradients:\\
      $\nabla_{\boldsymbol{\theta}} = \frac{1}{n}\mathbf{X}^\top(\hat{\boldsymbol{p}}-\boldsymbol{y})$\\
      $\nabla_{\theta_0} = \frac{1}{n}\sum_i(\hat{p}_i-y_i)$
    \item Update parameters:\\
      $\boldsymbol{\theta} \leftarrow \boldsymbol{\theta}-\alpha\,\nabla_{\boldsymbol{\theta}}$
  \end{enumerate}
\end{enumerate}
\end{block}
$\alpha$ = learning rate (hyperparameter)

\column{0.44\textwidth}
\begin{exampleblock}{Why It Works}
\begin{itemize}
  \item Loss is \textbf{convex} $\Rightarrow$ guaranteed convergence to global minimum
  \item No closed-form solution (unlike linear regression)
  \item Iterative optimization is necessary
\end{itemize}
\end{exampleblock}

\vspace{0.4em}
\begin{block}{Practical Variants}
\begin{itemize}
  \item \textbf{Batch GD}: all $m$ samples per step
  \item \textbf{Stochastic GD}: 1 random sample
  \item \textbf{Mini-batch GD}: $k$ samples ($k\ll m$)
\end{itemize}
\end{block}
\end{columns}
\end{frame}

% ------------------------------------------------------------
\begin{frame}{Linear vs Logistic rergression}
\begin{columns}[T]
\column{0.47\textwidth}
\begin{block}{Linear Regression}
\textbf{Loss:} $J = \|\boldsymbol{y}-\mathbf{X}\boldsymbol{\theta}\|^2$

\vspace{0.3em}
Set $\nabla J=\boldsymbol{0}$:
\[
  \mathbf{X}^\top\mathbf{X}\,\boldsymbol{\theta} = \mathbf{X}^\top\boldsymbol{y}
\]

\textbf{Closed-form solution:}
\[
  \hat{\boldsymbol{\theta}} = (\mathbf{X}^\top\mathbf{X})^{-1}\mathbf{X}^\top\boldsymbol{y}
\]
\end{block}

\column{0.47\textwidth}
\begin{block}{Logistic Regression}
\textbf{Loss:} $J = -\sum\!\bigl[y\log\sigma+(1{-}y)\log(1{-}\sigma)\bigr]$

\vspace{0.3em}
Set $\nabla J=\boldsymbol{0}$:
\[
  \mathbf{X}^\top\bigl(\sigma(\mathbf{X}\boldsymbol{\theta}+\theta_0\,\boldsymbol{1})-\boldsymbol{y}\bigr) = \boldsymbol{0}
\]

\textbf{\color{red!70} $\sigma(\cdot)$ is nonlinear $\Rightarrow$ cannot isolate $\boldsymbol{\theta}$!}
\end{block}
\end{columns}

\vspace{0.5em}
\centering
\textit{But the loss \textbf{is} convex $\Rightarrow$ iterative methods will find the \textbf{global optimum}.}
\end{frame}

% ============================================================
\section{Decision Boundaries \& Geometry}
% ============================================================
\begin{frame}{The Decision Boundary}
\begin{columns}[T]
\column{0.48\textwidth}
\begin{center}
\begin{tikzpicture}[scale=0.8]
  % Boundary goes from (-3,-2.5) to (4,3.0)
  % Slope = 5.5/7 ~ 0.786,  intercept: y = 0.786*x - 0.14
  % Red region: above the boundary line (class 0)
  % Blue region: below the boundary line (class 1)

  % Clipping rectangle
  \clip (-3,-3) rectangle (4.3,3.5);

  % Red shaded region (above boundary)
  \fill[red!6] (-3,-2.5) -- (4,3.0) -- (4,3.5) -- (-3,3.5) -- cycle;
  % Blue shaded region (below boundary)
  \fill[BrandBlue!6] (-3,-2.5) -- (4,3.0) -- (4,-3) -- (-3,-3) -- cycle;

  % Decision boundary
  \draw[dashed,BrandBlue,very thick] (-3,-2.5) -- (4,3.0);

  % Axes
  \draw[->,thick] (-3,0) -- (4.3,0) node[right,font=\small]{$x_1$};
  \draw[->,thick] (0,-3) -- (0,3.5) node[above,font=\small]{$x_2$};

  % Class 0 points — all above the line y = 0.786x - 0.14
  \foreach \x/\y in {-2.0/1.8, -1.4/1.5, -0.8/2.5, -1.8/1.0,
                     -0.3/2.0, -2.5/1.2, -1.0/1.8, 0.3/2.4,
                     -0.6/1.3, -1.5/2.6}{
    \fill[red!70] (\x,\y) circle (2.5pt);
  }
  % Class 1 points — all below the line y = 0.786x - 0.14
  \foreach \x/\y in {1.5/-0.8, 2.2/-0.5, 0.8/-1.5, 3.0/0.5,
                     2.5/-1.0, 1.8/-0.2, 1.0/-0.8, 3.2/-0.5,
                     2.0/0.8, 0.5/-1.8}{
    \fill[BrandBlue!80] (\x,\y) circle (2.5pt);
  }

  % Normal vector: perpendicular to boundary, pointing into class-1 region
  % Boundary direction (7, 5.5), normal toward class 1: (5.5, -7)
  % Normalised and scaled to length 1.3: (0.81, -1.03)
  \draw[->,very thick,accentorange] (0.5,0.25) -- (1.31,-0.78);
  \node[font=\scriptsize,accentorange] at (1.85,-0.55) {$\boldsymbol{\theta}$};

  % Boundary label
  \node[font=\scriptsize,BrandBlue,rotate=38] at (2.8,2.8) {$\boldsymbol{\theta}^\top\boldsymbol{x}+\theta_0=0$};

  \node[font=\scriptsize,red!70] at (-2.2,3.2) {$\hat{y}=0$};
  \node[font=\scriptsize,BrandBlue] at (3.5,-2.2) {$\hat{y}=1$};
\end{tikzpicture}
\end{center}

\column{0.48\textwidth}
\textbf{The decision boundary is defined by:}
\[
  \boldsymbol{\theta}^\top\boldsymbol{x}+\theta_0=0
\]

This is a \textbf{hyperplane} in feature space:
\begin{itemize}
  \item 2D features: a line
  \item 3D features: a plane
  \item $p$D features: a hyperplane
\end{itemize}

\vspace{0.3em}
\textbf{Geometric properties:}
\begin{itemize}
  \item $\boldsymbol{\theta}$ is \textbf{normal} (perpendicular) to the boundary
  \item On the boundary: $P(y{=}1)=\sigma(0)=0.5$
  \item $\boldsymbol{\theta}^\top\boldsymbol{x}+\theta_0 > 0 \;\Rightarrow\;$ predict class 1
  \item $\boldsymbol{\theta}^\top\boldsymbol{x}+\theta_0 < 0 \;\Rightarrow\;$ predict class 0
\end{itemize}
\end{columns}
\end{frame}

% ------------------------------------------------------------
\begin{frame}{Probability Landscape}
\begin{columns}[T]
\column{0.50\textwidth}
\begin{center}
\begin{tikzpicture}[scale=0.75]
  % Previous boundary: y = 0.786*x - 0.14, i.e. direction (7, 5.5)
  % Normal toward class 1 (below-right): (5.5, -7), normalised: (0.617, -0.787)
  % So theta ~ (0.617, -0.787), and z = 0.617*x1 - 0.787*x2
  % Decision boundary: 0.617*x1 - 0.787*x2 = 0  =>  x2 = 0.784*x1 (positive slope, matching previous)

  % Draw a heatmap-style representation of probability
  \foreach \i in {-4,-3.5,...,4}{
    \foreach \j in {-3,-2.5,...,3}{
      \pgfmathsetmacro{\zval}{0.617*\i - 0.787*\j}
      \pgfmathsetmacro{\pval}{1/(1+exp(-2.0*\zval))}
      \pgfmathsetmacro{\rcomp}{100*(1-\pval)}
      \pgfmathsetmacro{\bcomp}{100*\pval}
      \fill[red!\rcomp!BrandBlue!\bcomp, opacity=0.5] (\i-0.25,\j-0.25) rectangle (\i+0.25,\j+0.25);
    }
  }

  % Decision boundary: x2 = 0.784*x1 (same positive slope as previous slide)
  \draw[white,very thick,dashed] (-3.5,-2.74) -- (4,3.14);

  % Axes
  \draw[->,thick] (-4.3,0) -- (4.5,0) node[right,font=\small]{$x_1$};
  \draw[->,thick] (0,-3.3) -- (0,3.5) node[above,font=\small]{$x_2$};

  % Contour labels — red (class 0) is above-left, blue (class 1) is below-right
  \node[font=\tiny,white,fill=black!50,rounded corners=1pt,inner sep=1pt] at (-1.8,2.2) {$p=0.1$};
  \node[font=\tiny,white,fill=black!50,rounded corners=1pt,inner sep=1pt] at (0.2,-0.3) {$p=0.5$};
  \node[font=\tiny,white,fill=black!50,rounded corners=1pt,inner sep=1pt] at (2.5,-1.5) {$p=0.9$};

  % Normal vector: theta ~ (0.617, -0.787), scaled to visible length
  % Points into class-1 region (below-right), consistent with previous slide
  \draw[->,very thick,accentorange] (0,0) -- (0.93,-1.18);
  \node[font=\scriptsize,accentorange,right] at (0.93,-1.18) {$\boldsymbol{\theta}$};
\end{tikzpicture}
\end{center}

\column{0.46\textwidth}
\textbf{The sigmoid creates a smooth gradient:}
\begin{itemize}
  \item Probability transitions smoothly from 0 to 1
  \item $\boldsymbol{\theta}$ points toward the \textbf{increasing-probability} direction
  \item $\|\boldsymbol{\theta}\|$ controls the \textbf{steepness} of the transition
\end{itemize}

\begin{block}{Interpretation of $\|\boldsymbol{\theta}\|$}
\begin{itemize}
  \item Large $\|\boldsymbol{\theta}\|$: sharp transition, the model is \textbf{very confident}
  \item Small $\|\boldsymbol{\theta}\|$: gradual transition, model is \textbf{uncertain}
  \item Regularization controls $\|\boldsymbol{\theta}\|$ and thus model confidence
\end{itemize}
\end{block}
\end{columns}
\end{frame}

% ------------------------------------------------------------
\begin{frame}{Nonlinear Decision Boundaries via Feature Engineering}
\begin{columns}[T]
\column{0.40\textwidth}
\begin{center}
\vspace{-2em}
\textbf{Original feature space} $(x_1, x_2)$\\[0.3em]
\begin{tikzpicture}[scale=0.75]
  % Axes
  \draw[->,thick] (-3.3,0) -- (3.5,0) node[right,font=\small]{$x_1$};
  \draw[->,thick] (0,-3.3) -- (0,3.5) node[above,font=\small]{$x_2$};

  % Inner class (class 0, red) — ring of points near origin
  \foreach \a in {0,35,...,325}{
    \pgfmathsetmacro{\r}{0.6+0.4*rnd}
    \fill[red!70] ({\r*cos(\a)},{\r*sin(\a)}) circle (2.5pt);
  }

  % Outer class (class 1, blue) — ring of points further out
  \foreach \a in {10,30,...,350}{
    \pgfmathsetmacro{\r}{1.8+0.7*rnd}
    \fill[BrandBlue!80] ({\r*cos(\a)},{\r*sin(\a)}) circle (2.5pt);
  }

  % Show that a linear boundary fails
  \draw[dashed,gray,thick] (-2.8,2.8) -- (2.8,-2.8);
  \node[font=\scriptsize,gray,rotate=-45] at (2.3,-1.8) {linear fails!};

  % Nonlinear (circular) decision boundary
  \draw[very thick,accentorange,dashed] (0,0) circle (1.35);
  \node[font=\scriptsize,accentorange] at (0,-1.5) {$\boldsymbol{\theta}^\top\boldsymbol{\phi}(\boldsymbol{x})+\theta_0=0$};

  % Labels
  \node[font=\scriptsize,red!70] at (0,0.35) {Class 0};
  \node[font=\scriptsize,BrandBlue] at (2.3,2.5) {Class 1};
\end{tikzpicture}

\vspace{0.3em}
{\footnotesize After mapping $\boldsymbol{x}\mapsto\boldsymbol{\phi}(\boldsymbol{x})$, the boundary\\
is \textbf{linear} in the lifted space but\\
\textbf{nonlinear} in the original space.}
\end{center}

\column{0.6\textwidth}
\vspace{-0.5em}
Logistic regression is \textbf{linear} in feature space, but we can create nonlinear boundaries by \textbf{lifting} the features:

\vspace{-0.5em}
\begin{block}{Nonlinear expansions}
\vspace{-0.5em}
\[
  \boldsymbol{\phi}(\boldsymbol{x}) = [x_1,\, x_2,\, x_1^2,\, x_1 x_2,\, x_2^2]^\top
\]
\vspace{-0.8em}
\begin{itemize}
  \item Degree 2 $\Rightarrow$ \textbf{conic sections} (circles, ellipses, hyperbolas)
  \item Degree $d$ $\Rightarrow$ $\binom{p+d}{d}$ features
  \item Risk of overfitting at high $d$--- use \textbf{regularization}!
\end{itemize}
\end{block}
\begin{alertblock}{Linear transforms $\neq$ nonlinearity}
PCA scores, PLSR latent variables, and ICA components are \textbf{linear} projections of $\mathbf{X}$. Using them as inputs still yields a \textbf{linear} boundary in the original space. 
\end{alertblock}
\end{columns}
\end{frame}

% ============================================================
\section{Model Evaluation}
% ============================================================

\begin{frame}{Confusion Matrix}
\begin{columns}[T]
\column{0.45\textwidth}
\begin{center}
\begin{tikzpicture}[scale=0.75,font=\small]
  % Headers
  \node[font=\normalsize\bfseries] at (2.5,4.8) {Predicted};
  \node[font=\small] at (1.25,4.2) {Positive};
  \node[font=\small] at (3.75,4.2) {Negative};
  \node[font=\normalsize\bfseries,rotate=90] at (-1.2,2) {Actual};
  \node[font=\small,rotate=90] at (-0.4,3.0) {Positive};
  \node[font=\small,rotate=90] at (-0.4,1.0) {Negative};
  
  % TP cell
  \fill[BrandBlue!15,rounded corners=3pt] (0,2.1) rectangle (2.4,3.9);
  \draw[BrandBlue,thick,rounded corners=3pt] (0,2.1) rectangle (2.4,3.9);
  \node[font=\Large\bfseries,BrandBlue] at (1.2,3.3) {TP};
  \node[font=\scriptsize,BrandBlue] at (1.2,2.6) {True Positive};
  
  % FP cell
  \fill[red!10,rounded corners=3pt] (2.6,2.1) rectangle (5.0,3.9);
  \draw[red!60,thick,rounded corners=3pt] (2.6,2.1) rectangle (5.0,3.9);
  \node[font=\Large\bfseries,red!70] at (3.8,3.3) {FP};
  \node[font=\scriptsize,red!70] at (3.8,2.6) {False Positive};
  
  % FN cell
  \fill[red!10,rounded corners=3pt] (0,0) rectangle (2.4,1.9);
  \draw[red!60,thick,rounded corners=3pt] (0,0) rectangle (2.4,1.9);
  \node[font=\Large\bfseries,red!70] at (1.2,1.3) {FN};
  \node[font=\scriptsize,red!70] at (1.2,0.5) {False Negative};
  
  % TN cell
  \fill[BrandBlue!15,rounded corners=3pt] (2.6,0) rectangle (5.0,1.9);
  \draw[BrandBlue,thick,rounded corners=3pt] (2.6,0) rectangle (5.0,1.9);
  \node[font=\Large\bfseries,BrandBlue] at (3.8,1.3) {TN};
  \node[font=\scriptsize,BrandBlue] at (3.8,0.5) {True Negative};
\end{tikzpicture}
\end{center}

\column{0.50\textwidth}
\textbf{Key Metrics:}

\vspace{0.3em}
$\text{Accuracy} = \dfrac{TP+TN}{TP+TN+FP+FN}$\\[0.2em]
{\scriptsize\textit{(misleading for imbalanced data!)}}

\vspace{0.4em}
$\text{Precision} = \dfrac{TP}{TP+FP}$\\[0.2em]
{\scriptsize\textit{Of predicted positives, how many truly are?}}

\vspace{0.4em}
$\text{Recall (Sensitivity)} = \dfrac{TP}{TP+FN}$\\[0.2em]
{\scriptsize\textit{Of all true positives, how many did we catch?}}

\vspace{0.4em}
$\text{F1} = 2\cdot\dfrac{\text{Precision}\cdot\text{Recall}}{\text{Precision}+\text{Recall}}$\\[0.2em]
{\scriptsize\textit{Harmonic mean---balances precision \& recall}}

\vspace{0.4em}
$\text{Specificity} = \dfrac{TN}{TN+FP}$
\end{columns}
\end{frame}

% ------------------------------------------------------------
\begin{frame}{ROC Curve and AUC}
\begin{columns}[T]
\column{0.45\textwidth}
\begin{center}
\begin{tikzpicture}[scale=0.7]
  % Axes
  \draw[->,thick] (0,0) -- (5.5,0) node[right,font=\small]{FPR};
  \draw[->,thick] (0,0) -- (0,5.5) node[above,font=\small]{TPR};
  \node[left,font=\tiny] at (0,0) {$0$};
  \node[left,font=\tiny] at (0,5) {$1$};
  \node[below,font=\tiny] at (5,0) {$1$};
  
  % AUC shading
  \fill[BrandBlue!10] (0,0) 
    plot[domain=0:5,samples=50] (\x, {5*(\x/5)^0.35}) -- (5,0) -- cycle;
  
  % Diagonal (random)
  \draw[dashed,gray] (0,0) -- (5,5);
  \node[font=\tiny,gray,rotate=45] at (3.8,4.2) {Random ($\text{AUC}=0.5$)};
  
  % ROC curve
  \draw[very thick,BrandBlue] plot[domain=0:5,samples=50] (\x, {5*(\x/5)^0.35});
  
  % AUC label
  \node[font=\normalsize\bfseries,BrandBlue] at (3.2,2.0) {AUC $= 0.91$};
\end{tikzpicture}
\end{center}

\column{0.50\textwidth}
\textbf{ROC} = Receiver Operating Characteristic

\vspace{0.3em}
\textbf{How it works:}
\begin{enumerate}
  \item Vary the classification threshold from $0$ to $1$
  \item At each threshold, compute TPR and FPR
  \item Plot TPR ($y$-axis) vs.\ FPR ($x$-axis)
\end{enumerate}

\vspace{0.3em}
\textbf{Interpreting AUC:}
\begin{itemize}
  \item $\text{AUC}=1.0$: perfect classifier
  \item $\text{AUC}=0.5$: random guessing (diagonal)
  \item $\text{AUC}<0.5$: worse than random
\end{itemize}

\vspace{0.3em}
\begin{exampleblock}{}
AUC is \textbf{threshold-independent}---excellent for comparing models across different operating points.
\end{exampleblock}
\end{columns}
\end{frame}

% ============================================================
\section{Regularization}
% ============================================================

\begin{frame}{Regularization}
\textbf{General form:}
\[
  J_{\text{reg}}(\boldsymbol{\theta},\theta_0) = J(\boldsymbol{\theta},\theta_0) + \lambda\,R(\boldsymbol{\theta}) \qquad\text{where }\lambda>0
\]

\begin{columns}[T]
\column{0.47\textwidth}
\begin{block}{L2 Regularization (Ridge)}
\[
  R(\boldsymbol{\theta}) = \|\boldsymbol{\theta}\|_2^2 = \sum_j \theta_j^2
\]
\vspace{-0.5em}
\begin{itemize}
  \item Shrinks all weights toward zero
  \item Distributes weight among correlated features
\end{itemize}
\end{block}

\column{0.47\textwidth}
\begin{block}{L1 Regularization (Lasso)}
\[
  R(\boldsymbol{\theta}) = \|\boldsymbol{\theta}\|_1 = \sum_j |\theta_j|
\]
\begin{itemize}
  \item Drives some weights \textbf{exactly to zero}
  \item Automatic \textbf{feature selection}
\end{itemize}
\end{block}
\end{columns}

\begin{center}
\begin{tikzpicture}[scale=0.55]
  % L2 contour
  \begin{scope}[shift={(-5,0)}]
    \draw[BrandBlue,thick] (0,0) circle (1.5);
    \fill[BrandBlue!10] (0,0) circle (1.5);
    \draw[->,thin] (-2.2,0) -- (2.2,0) node[right,font=\tiny]{$\theta_1$};
    \draw[->,thin] (0,-2.2) -- (0,2.2) node[above,font=\tiny]{$\theta_2$};
    \node[font=\scriptsize\bfseries,BrandBlue] at (0,-2.7) {L2 (circle)};
  \end{scope}
  % L1 contour
  \begin{scope}[shift={(3,0)}]
    \fill[accentorange!10] (0,1.5) -- (1.5,0) -- (0,-1.5) -- (-1.5,0) -- cycle;
    \draw[accentorange,thick] (0,1.5) -- (1.5,0) -- (0,-1.5) -- (-1.5,0) -- cycle;
    \draw[->,thin] (-2.2,0) -- (2.2,0) node[right,font=\tiny]{$\theta_1$};
    \draw[->,thin] (0,-2.2) -- (0,2.2) node[above,font=\tiny]{$\theta_2$};
    \node[font=\scriptsize\bfseries,accentorange] at (0,-2.7) {L1 (diamond)};
    \node[font=\tiny,accentorange] at (2.8,1.0) {Corners $\Rightarrow$};
    \node[font=\tiny,accentorange] at (2.8,0.5) {sparse solutions};
  \end{scope}
\end{tikzpicture}
\end{center}
\end{frame}

% ============================================================
\section{Extensions \& Wrap-up}
% ============================================================
\begin{frame}{Multiclass Logistic Regression}
\begin{columns}[T]
\column{0.45\textwidth}
\begin{center}
\begin{tikzpicture}[scale=0.65]
  % Three class clusters
  \fill[BrandBlue!8] (-1,2) circle (1.8);
  \fill[red!8] (3,2) circle (1.8);
  \fill[accentgreen!8] (1,-1.5) circle (1.8);
  
  \draw[BrandBlue!40,dashed] (-1,2) circle (1.8);
  \draw[red!40,dashed] (3,2) circle (1.8);
  \draw[accentgreen!40,dashed] (1,-1.5) circle (1.8);
  
  % Points class A
  \foreach \a in {0,40,...,320}{
    \pgfmathsetmacro{\r}{0.3+rnd*1.0}
    \fill[BrandBlue!80] ({-1+\r*cos(\a)},{2+\r*sin(\a)}) circle (2pt);
  }
  % Points class B
  \foreach \a in {0,40,...,320}{
    \pgfmathsetmacro{\r}{0.3+rnd*1.0}
    \fill[red!70] ({3+\r*cos(\a)},{2+\r*sin(\a)}) circle (2pt);
  }
  % Points class C
  \foreach \a in {0,40,...,320}{
    \pgfmathsetmacro{\r}{0.3+rnd*1.0}
    \fill[accentgreen!80] ({1+\r*cos(\a)},{-1.5+\r*sin(\a)}) circle (2pt);
  }
  
  \node[font=\small\bfseries,BrandBlue] at (-1,2) {A};
  \node[font=\small\bfseries,red!70] at (3,2) {B};
  \node[font=\small\bfseries,accentgreen] at (1,-1.5) {C};
\end{tikzpicture}
\end{center}

\column{0.50\textwidth}
\vspace{-1.0em}
\begin{block}{One-vs-Rest (OvR)}
\begin{itemize}
  \item Train $K$ separate binary classifiers
  \item Classifier $k$: class $k$ vs.\ all others
  \item Predict class with highest confidence
\end{itemize}
\end{block}
\vspace{-0.75em}
\begin{exampleblock}{Softmax (Multinomial)}
Direct generalization to $K$ classes:
\[
  P(y{=}k\mid\boldsymbol{x}) = \frac{e^{\boldsymbol{\theta}_k^\top\boldsymbol{x}}}{\sum_{j=1}^{K}e^{\boldsymbol{\theta}_j^\top\boldsymbol{x}}}
\]
\vspace{-1.0em}
\begin{itemize}
  \item All $K$ probabilities sum to 1
  \item Loss: categorical cross-entropy
  \item Preferred when classes are mutually exclusive
\end{itemize}
\end{exampleblock}
\end{columns}
\end{frame}

% ------------------------------------------------------------
\begin{frame}{Practical Considerations}
\begin{columns}[T]
\column{0.47\textwidth}
\begin{block}{Feature Scaling}
Standardize features (zero mean, unit variance) before fitting. Gradient descent converges much faster with scaled features.
\end{block}

\vspace{0.3em}
\begin{block}{Handling Imbalanced Data}
\begin{itemize}
  \item Use \texttt{class\_weight='balanced'}
  \item Consider oversampling (SMOTE)
  \item Adjust decision threshold via ROC
\end{itemize}
\end{block}

\column{0.47\textwidth}
\begin{block}{Feature Multicollinearity}
\begin{itemize}
  \item L2 regularization handles moderate multicollinearity
  \item For severe cases: use PCA or PLSR first to reduce dimensions
\end{itemize}
\end{block}

\vspace{0.3em}
\begin{alertblock}{Convergence Issues}
If classes are \textbf{perfectly separable}, $\boldsymbol{\theta}\to\pm\infty$. Regularization prevents this. Always check convergence warnings!
\end{alertblock}
\end{columns}
\end{frame}

% ------------------------------------------------------------
\begin{frame}{Where Does Logistic Regression Fit?}
\centering
\small
\begin{tabular}{llcll}
\toprule
\textbf{Method} & \textbf{Type} & \textbf{Supervised?} & \textbf{Outputs} & \textbf{Key Idea} \\
\midrule
Linear Reg. & Regression & Yes & Continuous $\hat{y}$ & Min.\ MSE \\
\textbf{Logistic Reg.} & \textbf{Classification} & \textbf{Yes} & \textbf{Probabilities} & \textbf{Max.\ likelihood} \\
PCA & Dim.\ Reduction & No & Components & Max.\ variance \\
PLSR & Regression & Yes & Continuous $\hat{y}$ & Max.\ covariance \\
t-SNE & Visualization & No & 2D/3D map & Preserve local structure \\
ICA & Source Sep. & No & Indep.\ sources & Max.\ independence \\
\bottomrule
\end{tabular}

\vspace{1em}
\begin{tikzpicture}
\node[draw=BrandBlue, fill=SoftBlue, rounded corners=4pt, inner sep=6pt, font=\footnotesize\bfseries, align=center] {%
  Pipeline: \; Data $\to$ PCA/PLSR $\to$ Logistic Regression $\to$ ROC/AUC%
};
\end{tikzpicture}
\end{frame}

% ------------------------------------------------------------
\begin{frame}{Key Takeaways}
\begin{enumerate}
  \item Logistic regression applies $\sigma(\boldsymbol{\theta}^\top\boldsymbol{x}+\theta_0)$ to map linear combinations to \textbf{probabilities}
  \item The loss function (binary cross-entropy) comes directly from \textbf{MLE with Bernoulli data}
  \item No closed-form solution, but \textbf{convexity} guarantees convergence of gradient descent
  \item The decision boundary $\boldsymbol{\theta}^\top\boldsymbol{x}+\theta_0=0$ is a \textbf{hyperplane}---linear in feature space
  \item \textbf{Regularization} (L1/L2) prevents overfitting and handles perfect separability
  \item Evaluation: use precision, recall, F1, \textbf{ROC/AUC}---not just accuracy
  \item Multiclass via One-vs-Rest or \textbf{Softmax}
  \item Feature engineering extends to nonlinear boundaries
\end{enumerate}

\end{frame}

% ============================================================
% THANK YOU
% ============================================================
\NTNUThankYou[adil.rasheed@ntnu.no]{Adil Rasheed}

\end{document}
